% =============================================================================
%  latex-agentic — multilingual / internationalization demo
%  Engine: XeLaTeX ONLY (compile with `latexmk` or `make pdf`).
%
%  One compilable document that shows SERIOUS i18n: Latin, Cyrillic, Greek,
%  right-to-left Arabic, and Chinese (CJK) — each with the right font and the
%  right language environment. It is written to DEGRADE GRACEFULLY: if a script
%  font is missing on your machine, that section prints a short note instead of
%  failing the build. See docs/INTERNATIONALIZATION.md for the full theory.
%
%  THE GOLDEN RULES OF XeLaTeX i18n (each illustrated below):
%    * No font covers every script. Assign a font PER script.
%    * Latin Modern (the LaTeX default) has NO Cyrillic and NO Greek — that is
%      why the Latin/Cyrillic/Greek font chain below matters.
%    * Arabic and Hebrew are right-to-left: polyglossia + the bidi engine
%      reorder the text. Load order matters (polyglossia pulls in bidi LAST).
%    * CJK needs xeCJK for correct line-breaking (no spaces between words).
% =============================================================================
\documentclass[11pt,a4paper]{article}
\usepackage{geometry}\geometry{margin=2.5cm}

% --- fontspec + polyglossia (XeLaTeX-native i18n) ---------------------------
\usepackage{fontspec}
\defaultfontfeatures{Ligatures=TeX}

\usepackage{polyglossia}
\setmainlanguage{english}
% Declare every OTHER language we use. \setotherlanguage gives us an
% environment (\begin{german}…) and the matching hyphenation patterns.
\setotherlanguage{german}
\setotherlanguage{french}
\setotherlanguage{russian}
\setotherlanguage{greek}
% Serbian is special: ONE language, TWO scripts. polyglossia selects the
% script with the [script=…] option (NOT variant=, which is the ekavian/
% ijekavian dialect switch) — see the Serbian section below.
\setotherlanguage[script=cyrillic]{serbian}
% Arabic is right-to-left; polyglossia loads the bidi package to handle it.
\setotherlanguage{arabic}

% ---------------------------------------------------------------------------
%  FONT STRATEGY (the heart of this template)
%
%  Latin / Cyrillic / Greek main font — try, in order:
%    1. "CMU Serif"      (Computer Modern Unicode; full Latin+Cyrillic+Greek)
%    2. "FreeSerif"      (GNU FreeFont; ships with TeX Live; very wide coverage)
%    3. Latin Modern Roman (ALWAYS present, but NO Cyrillic/Greek — last resort)
%
%  We test the family name first, then the TeX Live *filename* (which always
%  resolves inside the TeX tree even when fontconfig has not indexed it).
% ---------------------------------------------------------------------------
\IfFontExistsTF{CMU Serif}{%
  \setmainfont{CMU Serif}%
}{\IfFontExistsTF{cmunrm.otf}{%
  \setmainfont{CMU Serif}[%
    Extension=.otf, UprightFont=cmunrm, BoldFont=cmunbx,
    ItalicFont=cmunti, BoldItalicFont=cmunbi]%
}{\IfFontExistsTF{FreeSerif}{%
  \setmainfont{FreeSerif}%
}{\IfFontExistsTF{FreeSerif.otf}{%
  \setmainfont{FreeSerif}[%
    Extension=.otf, UprightFont=FreeSerif, BoldFont=FreeSerifBold,
    ItalicFont=FreeSerifItalic, BoldItalicFont=FreeSerifBoldItalic]%
}{%
  % LAST RESORT — note: Latin Modern lacks Cyrillic and Greek glyphs, so the
  % Russian/Greek/Serbian-Cyrillic sections may show "missing glyph" boxes.
  \setmainfont{Latin Modern Roman}%
}}}}

% Greek-aware font: CMU Serif and FreeSerif both cover Greek, so we reuse the
% main font for Greek. (polyglossia will use \greekfont if defined; if not, it
% falls back to the main font, which here already has Greek.)

% ---------------------------------------------------------------------------
%  Arabic font chain. Arabic NEEDS shaping (contextual joining), so the font
%  must be a real Arabic face. Try: Amiri → Scheherazade New → Noto Naskh
%  Arabic. If none is found we set a flag and SKIP the Arabic body (the doc
%  still compiles) — the section prints a short note instead.
% ---------------------------------------------------------------------------
\newif\ifhasarabic
\IfFontExistsTF{Amiri}{%
  \newfontfamily\arabicfont[Script=Arabic]{Amiri}\hasarabictrue
}{\IfFontExistsTF{Amiri-Regular.ttf}{%
  \newfontfamily\arabicfont[Script=Arabic]{Amiri-Regular.ttf}\hasarabictrue
}{\IfFontExistsTF{Scheherazade New}{%
  \newfontfamily\arabicfont[Script=Arabic]{Scheherazade New}\hasarabictrue
}{\IfFontExistsTF{Noto Naskh Arabic}{%
  \newfontfamily\arabicfont[Script=Arabic]{Noto Naskh Arabic}\hasarabictrue
}{%
  \hasarabicfalse
}}}}

% ---------------------------------------------------------------------------
%  Chinese (CJK) via xeCJK. Try FandolSong (ships with TeX Live) → Noto Serif
%  CJK SC → Songti SC. Same skip-with-note fallback if nothing is available.
%  IMPORTANT: load xeCJK and set the CJK font only when we have one, otherwise
%  \setCJKmainfont errors out.
% ---------------------------------------------------------------------------
\newif\ifhascjk
\IfFontExistsTF{FandolSong-Regular.otf}{\hascjktrue}{%
  \IfFontExistsTF{Noto Serif CJK SC}{\hascjktrue}{%
    \IfFontExistsTF{Songti SC}{\hascjktrue}{\hascjkfalse}}}

\ifhascjk
  \usepackage{xeCJK}
  \IfFontExistsTF{FandolSong-Regular.otf}{%
    % FandolSong-Regular has no built-in bold shape; pair it with the
    % matching bold file so \textbf works in CJK text.
    \setCJKmainfont{FandolSong-Regular.otf}[BoldFont=FandolSong-Bold.otf]%
  }{\IfFontExistsTF{Noto Serif CJK SC}{%
    \setCJKmainfont{Noto Serif CJK SC}%
  }{%
    \setCJKmainfont{Songti SC}%
  }}
\fi

% --- Typography & links (hyperref last) -------------------------------------
\usepackage{microtype}
\usepackage[autostyle=true]{csquotes}
\usepackage{xcolor}
\usepackage[
  unicode=true, hidelinks,
  pdftitle={Multilingual XeLaTeX demo},
  pdfauthor={latex-agentic}
]{hyperref}

\begin{document}

% --- Title block ------------------------------------------------------------
% NOTE (a real i18n trap): the standard \maketitle does not coexist with
% hyperref once a right-to-left language (here Arabic, via bidi) is active —
% the combination overflows TeX's input stack at \end{document} on this and
% similar TeX Live versions. We therefore typeset the title by hand. It looks
% the same, is easier to restyle, and sidesteps the bug entirely.
\begin{center}
  {\LARGE\bfseries A Multilingual XeLaTeX Document}\\[4pt]
  {\large latex-agentic}\\[2pt]
  {\today}
\end{center}
\vspace{1.2em}

\noindent This document is typeset with \textbf{XeLaTeX} and
\texttt{polyglossia}. Each section below is written in a different language and
script. The source is plain UTF-8 — the non-Latin text is typed directly, with
no escape codes. If a required font is not installed, the affected section
prints a short note rather than breaking the build.

% ===========================================================================
\section*{English (Latin, default)}
The quick brown fox jumps over the lazy dog. This paragraph uses the document's
default language and main serif font. Everything that follows switches language
and, where necessary, font.

% ===========================================================================
\section*{German (Latin)}
\begin{german}
Falsches Üben von Xylophonmusik quält jeden größeren Zwerg. Die deutsche
Sprache nutzt die Umlaute ä, ö, ü und das Eszett ß. Polyglossia lädt zudem die
korrekten Trennmuster für den Blocksatz.
\end{german}

% ===========================================================================
\section*{French (Latin)}
\begin{french}
Voix ambiguë d'un cœur qui, au zéphyr, préfère les jattes de kiwis. Le français
emploie de nombreux accents — é, è, ê, ë, à, ç — et des guillemets « à la
française ». La typographie française insère une espace fine avant ; : ? et !
\end{french}

% ===========================================================================
\section*{Serbian — dual script (Latin and Cyrillic)}
% Serbian is the textbook "one language, two scripts" case. We declared
% \setotherlanguage[script=cyrillic]{serbian} in the preamble; the [script=…]
% option switches scripts per block (variant= would switch the ekavian/
% ijekavian dialect instead). Below: first Latin, then Cyrillic.
\textbf{Latinica:}
\begin{serbian}[script=latin]
Đačko društvo čeka prijatan, žustar i živahan ćuk. Srpski se piše dvama
ravnopravnim pismima; ovo je latinica.
\end{serbian}

\smallskip\noindent\textbf{Ћирилица:}
\begin{serbian}[script=cyrillic]
Ђачко друштво чека пријатан, жустар и живахан ћук. Ово је исти текст исписан
ћирилицом — потпуно равноправним српским писмом.
\end{serbian}

% ===========================================================================
\section*{Russian (Cyrillic)}
\begin{russian}
Съешь же ещё этих мягких французских булок, да выпей чаю. Это известная
русская панграмма. Кириллический текст требует шрифта с поддержкой кириллицы —
Latin Modern такой поддержки не имеет.
\end{russian}

% ===========================================================================
\section*{Greek}
\begin{greek}
Ξεσκεπάζω την ψυχοφθόρα βδελυγμία. Η ελληνική γλώσσα χρησιμοποιεί το δικό της
αλφάβητο, με τόνους και διαλυτικά. Η γραμματοσειρά πρέπει να καλύπτει το
ελληνικό σύστημα γραφής.
\end{greek}

% ===========================================================================
\section*{Arabic (right-to-left)}
% Arabic flips the paragraph direction. The environment that switches BOTH the
% font and the base direction is the CAPITALISED language name —
% \begin{Arabic}…\end{Arabic}. (The lower-case \arabic / \textarabic forms are
% the *inline* commands.) Using the wrong form for RTL is a classic trap: the
% text silently falls back to the Latin main font — which has no Arabic glyphs
% — and the run can blow the input stack. We only enter the environment if an
% Arabic font was found (see the \ifhasarabic flag in the preamble).
\ifhasarabic
\begin{Arabic}
مرحبًا بالعالم. هذا نصٌّ تجريبيٌّ مكتوبٌ باللغة العربية، ويُكتب من اليمين إلى
اليسار. لاحظ كيف تتشابك الحروف وتتغيّر أشكالها حسب موضعها في الكلمة.
\end{Arabic}
\else
\textit{Arabic section skipped: no Arabic font found. Install one of
``Amiri'', ``Scheherazade New'', or ``Noto Naskh Arabic'' (e.g.
\texttt{tlmgr install amiri}) and rebuild.}
\fi

% ===========================================================================
\section*{Chinese (CJK)}
% CJK has no spaces between words, so line-breaking is handled glyph-by-glyph
% — that is xeCJK's job. We only typeset this if a CJK font was found.
\ifhascjk
你好，世界。这是一段用中文书写的示例文字。中文不使用空格来分隔词语，
因此换行需要 xeCJK 在字符之间进行处理。汉字与拉丁字母可以混排。
\else
\textit{Chinese section skipped: no CJK font found. Install one of
``FandolSong'', ``Noto Serif CJK SC'', or ``Songti SC'' (e.g.
\texttt{tlmgr install fandol}) and rebuild.}
\fi

\end{document}
